\justifying
\NSFCBodyText

申请人正在主持的、与本课题方法相关的项目如下。所列经费与编号以已有书面材料为准；未写入编号的项目不编造批准号。

（1）自治区天池英才引进计划（青年博士人才）项目“多模态融合技术研究及其在新疆产业集群中的应用”，新疆人社厅，经费 50 万元，2025.02--2028.02。申请人主持。该项目面向产业集群场景中的多模态融合，与本课题的维吾尔语场景文字及维—汉图像理解对象不同，不构成相同内容重复申报。

（2）自治区高校基本科研业务费科研项目（培育类）“融合多源数据大模型驱动的光伏发电功率预测方法研究”，项目编号 XJEDU2026P005，新疆教育厅，2026.01--2026.12。申请人主持。研究对象为光伏功率预测，仅提供深度学习训练经验，不构成本课题研究内容。

（3）自治区优秀博士后资助“面向通用领域的多模态视觉问答关键技术研究”，新疆人社厅，经费 3 万元，2024.07--2025.11。申请人主持。该项目提供视觉问答训练基础，已近结束；本课题将其方法迁移到民族语言场景图文，而不是重复通用域 VQA 任务。

（4）横向课题“面向政企数字化应用的智能知识库与问答系统研发”，项目编号 202604140013，经费 12 万元，2026.06--2026.10。申请人主持。该项目为应用系统研发，与本课题的公开基准评测无相同考核指标。

主要参与研究生的其他在研课题与本开放课题无相同研究内容。申请人目前未主持尚未结题的本实验室开放课题。
